\documentclass{article}
\usepackage{arxiv}
\usepackage[utf8]{inputenc}
\usepackage[T1]{fontenc}
\usepackage{booktabs}
\usepackage{hyperref}
\title{LinkerForge: An Auditable Prototype for Linker-Candidate Mining and Geometry-Guided Fusion Design}
\author{LinkerForge Research Team}
\date{19 October 2026}
\begin{document}
\maketitle
\begin{abstract}
We describe a computational prototype combining a sequence-based disordered flexible linker (DFL) classifier, a domain-gap candidate filter, provenance-stratified retrieval, geometric feasibility heuristics, attachment-site scoring, and pairwise ranking. A three-layer convolutional detector achieved reported mean residue average precision $0.7611\pm0.0110$ across three seeds on a family-disjoint development split, but this is not a proteome-mining precision estimate. A strict, exact-sequence-matched InterPro/Pfam audit found only one eligible gap in 28 development proteins, covering 38 of 547 known positive residues and no known negative residues. A separate builder reported 93 true-positive and one false-positive known residues under a different domain-annotation path. These evaluations cannot establish prospective span-level precision. The builder reports 27,595 model-scored gaps, which are candidates rather than verified DFLs. Retrieval, position-matched attachment-site, and cross-cluster real-decoy ranking results remain retrospective diagnostics. No independent no-linker specificity, published-method superiority, or experimental fusion performance has been demonstrated.
\end{abstract}

\section{Scope and evidence boundary}
Domain gaps, curated DFL annotations, structural domain boundaries, engineered linker templates, and model predictions are separate evidence types. A gap between domains is not necessarily disordered or flexible. Unannotated residues are unknown, not negative. CAID3 has already been used in a retrospective comparison and cannot be reused as a new blind test. The earlier version of this manuscript overstated calibration, dataset validation, and downstream performance; this version supersedes it. The detailed corrections are recorded in \texttt{reports/PAPER\_CLAIMS\_AUDIT\_2026-10-19.md}.

\section{Prototype methods}
The detector is a three-layer 1D CNN trained with a mask over observed positive and negative residues. A deterministic gate proposes 5--60-residue intervals between adjacent, non-overlapping, distinct domain annotations. Its code does not prove that the domains are folded or the interval is a DFL. The builder scores these intervals using a frozen CNN. A threshold routine selects among residue-score cutoffs on development data and computes one-sided Clopper--Pearson upper bounds at each cutoff. This is a descriptive calculation, not a valid split-conformal or learn-then-test guarantee: the same labels select the cutoff, there is no multiplicity correction, residues within proteins are correlated, and production span calls use a different mean/maximum-score rule. The nominal bound must not be extrapolated to new proteins or to entire mined spans.

The design module implements necessary length and geometry checks, retrieval from source-stratified records, an attachment-site heuristic, and a pairwise ranker. A geometric pass is not proof of a realizable structure. For attachment-site scoring, an explicit leakage diagnostic shows that the earlier terminal-vs-interior task is separated by window-edge distance alone (\texttt{term\_prox} AUROC $1.0000$, \texttt{rel\_interior} AUROC $0.9998$); under a position-matched 5-fold cluster-disjoint benchmark with positional weights fixed to zero, sequence-only features score near chance (AUROC $0.4756$ on natural authorities and $0.4589$ on mined candidates), demonstrating that attachment-site selection requires 3D structural features rather than 1D sequence propensities. For candidate ranking, we evaluate a zero-initialized pairwise ranker on cross-cluster real-decoy ranking (1 held-out true linker against 19 length- and rigidity-matched real linkers from other homology clusters) and non-oracle geometry-derived lengths on the 91 coordinate-complete CATH linkers.

\section{Recorded diagnostics}
\begin{table}[h]
\centering
\small
\begin{tabular}{p{0.39\linewidth}p{0.50\linewidth}}
\toprule
Diagnostic & Result and supported interpretation \\
\midrule
CNN development discrimination & Three-seed residue APS $0.7611\pm0.0110$ and AUROC $0.8583\pm0.0080$ on known labels; not mining precision. \\
Builder's gated development cutoff & 93 TP, 1 FP at residue cutoff $0.639120$ on known residues; descriptive, adaptively selected, and not a span-level guarantee. \\
Independent InterPro/Pfam audit & One eligible gap in 28 proteins; 38 known positive, zero known negative and three unknown residues in that gap. Precision cannot be estimated. \\
Control panel & 39/40 proteins without a high-tier call; the single/homo-domain subset fails the gate by construction and the panel lacks independent no-linker adjudication. \\
Reported inventory & 6,556 high-tier and 21,039 provisional model-scored gaps, plus 91 CATH-derived, 57 DisProt-curated and 23 template records. Model scores are not biological ground truth. \\
Retrieval and non-oracle CATH generators & 43/91 CATH queries within $\pm2$ length residues in the X2 hybrid arm; 37/91 ($40.66\%$) within $\pm2$ residues under non-oracle \texttt{design\_minimum\_plus\_2}. \\
Attachment-site position-matched test & Original task solved by \texttt{term\_prox} alone (AUROC $1.0000$); position-matched 5-fold CV yields AUROC $0.4756$ ($95\%$ CI $[0.4162, 0.5300]$) on natural and $0.4589$ on mined windows. \\
Cross-cluster real-decoy ranking & Natural ($148$ queries vs 19 length/rigidity-matched cross-cluster decoys): Hit@1 $0.2635$ ($95\%$ CI $[0.1875, 0.3355]$) vs $0.0500$ chance, MRR $0.4260$ vs $0.1799$ chance; Mined ($598$ queries): Hit@1 $0.1405$, MRR $0.2996$. Does not measure experimental fusion activity. \\
\bottomrule
\end{tabular}
\caption{Numbers transcribed from tracked reports; interpretation is limited to their actual evaluation tasks.}
\end{table}

The builder targeted a nominal residue-level FDR of 0.10, not 0.05. Its selected $0.639120$ cutoff is different from the production span rule, which combines a high-residue cutoff of $0.735198$ with a provisional mean cutoff of 0.50. The former draft's $0.92285$ threshold, 44.74\%-to-1.06\% proteome-wide FDR reduction, and 95\% precision guarantee are not supported by the tracked builder and summary.

The 184,554 unreviewed reference-proteome domain-gap observations are a separate stage-one audit, not 184,554 additional CNN-scored proteins. Existing fusion PDB exports (\texttt{reports/structures/sfgfp\_mcherry\_*}) are schematic coordinate assemblies on real crystal structures (\texttt{2B3P} and \texttt{2H5Q}), not experimentally resolved, Rosetta-validated, AlphaFold-validated, or molecular-dynamics-validated structures.

\section{Requirements before a performance claim}
The domain-annotation pipelines must be reconciled, and the exact production caller evaluated at span level on a sealed, family-disjoint set containing independently adjudicated eligible positives and no-linker negatives. Calibration must use an appropriate cluster-aware, prespecified procedure and account for threshold selection. Source data, weights, versions and hashes need reproducible release where licensing permits. Published predictors must be compared on identical targets, input modalities where possible, and residue coverage. No result is guaranteed in advance.

\section{Conclusion}
LinkerForge is an integrated computational prototype for generating and ranking linker candidates. Current evidence supports implementation and exploratory diagnostic claims, not certified high-precision proteome mining, general attachment-site accuracy, functional fusion design, or superiority over published methods.

\end{document}
